\documentclass[10pt,a4paper]{amsart}
\usepackage[margin=19mm]{geometry}
\usepackage{amsmath,amssymb,mathtools}
\usepackage[scr=rsfs,cal=euler]{mathalfa}
\usepackage{xfrac}
\usepackage[unicode,hidelinks,bookmarks=false]{hyperref}
\newcommand{\ie}{i.e.\ }
\newcommand{\cf}{cf.\ }
\newcommand{\resp}{respectively\ }
\newcommand{\Subsection}{\subsection}
\providecommand{\nobreakdash}{\nobreak\hbox{-}}
\setlength{\emergencystretch}{2em}
\multlinegap0pt
\usepackage{amssymb}
\usepackage{mathtools}
\usepackage[shortlabels]{enumitem}
\newtheorem{definition}{Definition}
\newtheorem{lemma}{Lemma}
\title{\large{\bfseries On the maximum number of triangles in tripartite graphs with no $4$-cycles between any two parts}}
\author{Chunqiu Fang, Rongxing Xu}
\date{Source: arXiv:2609.06594v1 (CC BY 4.0)}
\begin{document}
\maketitle

\noindent\emph{Source and licence.} \large{\bfseries On the maximum number of triangles in tripartite graphs with no $4$-cycles between any two parts}. Version arXiv:2609.06594v1 (CC BY 4.0), distributed by arXiv under the Creative Commons Attribution 4.0 International licence, http://creativecommons.org/licenses/by/4.0/, as recorded in the arXiv licence metadata for this exact version and shown on the abstract page https://arxiv.org/abs/2609.06594v1. Full licence notice: \texttt{licenses/arxiv-2609.06594-CC-BY-4.0.txt}.\par\smallskip

\noindent\textit{Editorial scope.} Counted unit: the article's lemma lem:planar-graph (source0.tex lines 228--234). The article's complete original proof of it is delivered with it (source0.tex lines 236--261, one \texttt{proof} environment of 26 lines). No mathematics is written, repaired or re-derived: the statement and the proof are the article's own lines, in the article's own notation, and no retained line is substituted or re-flowed.  Retained uncounted as the source-local context the proof needs: lines 165--167.  Nothing was omitted from any retained span, so the package carries no omission record.  No retained line contains a citation, so the wrapper carries no bibliography and no bibliography entry.  No retained line contains a figure environment, an image-inclusion command or drawing code, so no image is delivered and the package declares no asset dependency.  Editorial formatting only, in addition: an AMS \texttt{amsart} preamble that restates the article's own declarations and macros used by the retained lines, namely the environments \texttt{definition}, \texttt{lemma} on separate counters, together with the standard packages those lines need.  The retained environments print in this document as Definition 1, Lemma 1 in order of appearance, whereas the article prints the same items with the numbering of its own full document.  The complete native source of the article as distributed in the arXiv e-print for this exact version is bundled unchanged as \texttt{sources/209551/source0.tex}.
\begin{definition}\label{def:planar}
Let $F$ be a finite field and let $f\in F[X]$. The polynomial $f$ is \emph{planar over $F$} if, for every $a\in F^*=F\setminus\{0\}$, the map $x\mapsto f(x+a)-f(x)$ is a bijection from $F$ to itself.
\end{definition}

\begin{lemma}\label{lem:planar-graph}
Let $F$ be a finite field and let $f\in F[X]$ be planar over $F$. Let $A$ and $B$ be disjoint copies of $F^2$. For each $(x,y)\in F^2$, denote its copies in $A$ and $B$ by $(x,y)_A$ and $(x,y)_B$, respectively. Let $G$ be the bipartite graph with vertex set $A\cup B$ and edge set
\[
        E(G)=\bigl\{\{(x,y)_A,(x+t,y+f(t))_B\}:x,y,t\in F\bigr\}.
\]
Then $G$ does not contain a $4$-cycle.
\end{lemma}

\begin{proof}
Suppose to the contrary that two distinct vertices $u_1,u_2\in A$ have two distinct common neighbors $v_1,v_2\in B$. In the following coordinate calculations, we identify each vertex with its underlying element of $F^2$. For suitable $x_1,x_2,y_1,y_2\in F$,
\begin{align*}
v_1&=u_1+(x_1,f(x_1))=u_2+(y_1,f(y_1)),\\
v_2&=u_1+(x_2,f(x_2))=u_2+(y_2,f(y_2)).
\end{align*}
Thus we have
\[
 u_2-u_1
 =(x_1-y_1,f(x_1)-f(y_1))
 =(x_2-y_2,f(x_2)-f(y_2)).
\]
Comparing the first coordinates, we have
$x_1-y_1=x_2-y_2$. Let $t$ denote this common value.
Suppose that $t=0$. Then $x_1=y_1$ and $x_2=y_2$. In particular, the preceding vector equality becomes
\[
 u_2-u_1=(x_1-y_1,f(x_1)-f(y_1))=(0,0).
\]
It follows that $u_1=u_2$, contrary to the choice of two distinct vertices $u_1$ and $u_2$. Therefore, $t\ne0$.
Comparing the second coordinates, we obtain
\[
f(y_1+t)-f(y_1)=f(y_2+t)-f(y_2).
\]
By Definition~\ref{def:planar}, the map
$x\mapsto f(x+t)-f(x)$ is injective. Hence $y_1=y_2$. Since $x_1-y_1=x_2-y_2$, we also have $x_1=x_2$. It follows that $v_1=v_2$, a contradiction. Having two vertices in one part with two common neighbors is equivalent to having a $4$-cycle, so $G$ is $C_4$-free.
\end{proof}

\end{document}
